(a) A center $x$ gives a morphism $\operatorname{Spec}R\to X$ extending the canonical morphism $\operatorname{Spec}K\to X$, by Lemma 4.4. Separatedness makes this extension unique, so the image of the closed point, and hence the center, is unique.

(b) Properness gives the existence of that extension by the valuative criterion. Its closed-point image is a center, which is unique by (a).

(c) Suppose every valuation of $K/k$ has at most one center on $X$. Give the closure of the diagonal in $X\times_k X$ its reduced structure. It is integral with function field $K$. A valuation ring dominating the local ring at any of its points gives two morphisms $\operatorname{Spec}R\to X$ with the same generic map. Their centers coincide by hypothesis, so the morphisms coincide by Lemma 4.4. The chosen point therefore lies on the diagonal. The diagonal is an immersion with closed image, hence a closed immersion. Thus $X$ is separated.

Now suppose every valuation has exactly one center. The construction in \exref{II.4.10}, before the final use of properness, gives a proper surjective birational map $g:X'\to X$ with $X'$ integral and quasi-projective over $k$: it is closed in $X\times_k P$, where $P$ is projective, and closed in the open subset $P^0$ of $P$. Properness of $g$ lifts each center on $X$ uniquely to a center on $X'$. Embed $X'$ as an open subscheme of its integral projective closure. A boundary point would admit a valuation ring of $K$ dominating its local ring. That valuation would have two centers on the projective closure, one at the boundary and one in $X'$, contradicting (a). Hence $X'$ is projective. After any base extension, the image in the base of a closed subset of $X$ equals the image of its inverse image in $X'$, which is closed. Thus $X$ is universally closed over $k$, and therefore proper.

(d) Suppose $a\in\Gamma(X,\mathcal O_X)\setminus k$. Since $k$ is algebraically closed in every extension field, $a$ is transcendental over $k$. The local ring $k[a^{-1}]_{(a^{-1})}\subseteq K$ is dominated by a valuation ring $R$ of $K$, with $a^{-1}\in\mathfrak m_R$. By (b), $R$ has a center $x$ on $X$. Then $a\in\mathcal O_{X,x}\subseteq R$, contradicting $a^{-1}\in\mathfrak m_R$. Thus $\Gamma(X,\mathcal O_X)=k$.
